\subsubsection{The Spectrum of a Compact Operator}\label{subsec:compact-operator-spectrum}

In this section we prove a number of results about the spectrum of a compact operator. Since compact operators are bounded, the spectrum of a compact operator has all of the properties described in Section~\ref{subsec:bounded-operator-spectra}. Of course, with the added hypothesis of compactness, we can say a good bit more. We restrict ourselves to the case of operators on Hilbert space, though many of the results we give can be generalized to operators on Banach spaces.

In Hilbert spaces we can make use of the projection theorem and its consequences. In order to make use of this, we begin with a description of the spectrum of an operator of finite rank.

\begin{Lemma}
\label{lem:finite-rank-spectrum-alternative}
Suppose $A\in\mathcal{L}(H)$ has finite rank. Then for every $\lambda\in\mathbb{C}\backslash\{0\}$ exactly one of the following holds: either\begin{enumerate}\item $\lambda\in\rho(A)$, or\item $\lambda\in\sigma_p(A)$. In this case $\lambda$ is an eigenvalue of finite multiplicity.\end{enumerate}
\end{Lemma}

The proof follows directly from the corresponding result of linear algebra and is left to the reader (Problem~\ref{ex:finite-rank-spectrum-alternative-proof}).

We now prove a slightly different version of the Fredholm alternative theorem for operators of finite rank. This version is really just a technical result which will be useful in proving the analytic Fredholm theorem below (Theorem~\ref{thm:analytic-fredholm}).

\begin{Lemma}
\label{lem:analytic-finite-rank-alternative}
Let $G\subset\mathbb{C}$ be a domain, and suppose
\begin{equation}\label{eq:analytic-finite-rank-operator-family}
\mathbb{C}\supset G\ni\lambda\mapsto F(\lambda)\in\mathcal{L}(H)
\end{equation}
is analytic in $G$. Further suppose that, for every $\lambda\in G$, $F(\lambda)$ is of finite rank and that
\begin{equation}\label{eq:analytic-finite-rank-fixed-range-subspace}
\mathcal{R}(F(\lambda))\subseteq M,
\end{equation}
where $M$ is a finite-dimensional subspace of $H$, independent of $\lambda$. Then either\begin{enumerate}\item $(I-F(\lambda))^{-1}$ exists for no $\lambda\in G$, or\item $(I-F(\lambda))^{-1}$ exists for every $\lambda\in G\backslash S$ where $S$ is a discrete set in $G$ (i.e., it has no limit point in $G$). In this case the function $\lambda\mapsto(I-F(\lambda))^{-1}$ is analytic on $G\backslash S$, and if $\lambda\in S$, then $F(\lambda)\phi=\phi$ has a finite-dimensional family of solutions.\end{enumerate}
\end{Lemma}

\begingroup
\postdisplaypenalty=10000
\renewcommand{\endProof}{\par\nobreak\hfill$\blacksquare$\par\vspace{0.1cm}}
\begin{Proof}
Let $\{\psi_i\}_{i=1}^N$ be a basis for $M$. Then there are analytic vector functions
\begin{equation}\label{eq:analytic-finite-rank-vector-functions}
G\ni\lambda\mapsto\gamma_i(\lambda)\in H,\quad i=1,\ldots,N,
\end{equation}
such that
\begin{equation}\label{eq:analytic-finite-rank-family-expansion}
F(\lambda)\phi=\sum_{i=1}^n(\gamma_i(\lambda),\phi)\psi_i.
\end{equation}
Let $\Lambda(\lambda)$ be the $N\times N$ matrix with components
\begin{equation}\label{eq:analytic-finite-rank-matrix-components}
\Lambda_{ij}(\lambda)=(\gamma_j(\lambda),\psi_i).
\end{equation}
The reader should verify that $F(\lambda)\phi=\phi$ has a nontrivial solution if and only if
\begin{equation}\label{eq:analytic-finite-rank-invertibility-determinant}
d(\lambda):=\det(I-\Lambda(\lambda))=0.
\end{equation}
However, $d(\lambda)$ is analytic on $G$. Hence, by a standard result of complex variables, either $d$ is identically zero in $G$, or the zeros of $d$ form a discrete set.

Since the range of $F$ is finite-dimensional, so is the solution space of $F(\lambda)\phi=\phi$. This completes the proof.
\end{Proof}
\endgroup
